%==============================================================================
% Chapter 9 — End-to-End Systems Design & MLOps (Capstone)
% Curriculum: Phase 3, Chapter 9 of docs/AI_ML_Master_Checklist.md
% Companion code: projects/09_capstone/
%
% Section skeleton only. Figures for this chapter go in theory/figures/.
%==============================================================================
\chapter{End-to-End Systems Design \& MLOps}
\label{ch:capstone}

The capstone, and the only chapter that introduces no new model.\projectref{projects/09\_capstone/} It takes one
already trained in Phase~2 and does the work that turns it into something other
people can call: an HTTP interface, a durable log of what it was asked and what
it answered, and a container stack that starts with one command.

This is where the Chapter~\ref{ch:engineering-foundations} investment pays off.
The dependencies are already locked, the data already versioned, the pipeline
already runs on every pull request --- so the work here is genuinely about
service design rather than about repairing eight chapters of accumulated
improvisation. The distinction that organises the chapter is the training/serving
boundary: training is a batch job that may take hours and is allowed to fail,
while serving is a long-lived process that must answer in milliseconds and must
not.

%==============================================================================
\section{From Notebooks to Live Services}
\label{sec:notebooks-to-services}

\subsection{Why a Notebook Is Not a Deployment Artifact}

\subsection{The Training/Serving Boundary}

\subsection{Model Artifacts and Versioning}

\subsection{Training--Serving Skew}

%==============================================================================
\section{Data Persistence for Services}
\label{sec:persistence}

% No relational database in this chapter: the Compendium stays Python end to
% end, so durable state is an append-only log plus columnar files rather than
% Postgres. The design questions below are the ones that matter regardless of
% the storage engine.

\subsection{Durable versus Ephemeral State}

\subsection{Append-Only Event Logs}

\subsection{Schema Evolution for Logged Records}

\subsection{Why the Serving Path Must Not Block on a Write}

%==============================================================================
\section{Microservice Architecture}
\label{sec:microservices}

\subsection{Service Decomposition}

\subsection{Statelessness}

\subsection{Synchronous versus Asynchronous Boundaries}

\subsection{Health Checks}

\subsection{Twelve-Factor Configuration}

%==============================================================================
\section{RESTful APIs with FastAPI}
\label{sec:fastapi}

\subsection{REST Principles}

\subsection{Path and Query Parameters}

\subsection{Pydantic Request and Response Schemas}

\subsection{Dependency Injection}

\subsection{Wrapping a Model Behind \texttt{/predict}}

\subsection{OpenAPI Documentation}

%==============================================================================
\section{Prediction Logging in Python}
\label{sec:prediction-logging}

\subsection{Append-Only JSONL for Request/Response Records}

\subsection{Parquet via \texttt{pandas} and \texttt{pyarrow}}

\subsection{Partitioning by Date and Log Rotation}

\subsection{Reading the Log Back to Compute Drift}

% Closes the loop with \cref{sec:monitoring}: the logged records are the only
% evidence available that the live input distribution still resembles the
% training distribution.

%==============================================================================
\section{Containerization}
\label{sec:docker}

\subsection{Images versus Containers}

\subsection{Layer Caching}

\subsection{Multi-Stage Builds for Slim Python Images}

%==============================================================================
\section{Orchestration with Docker Compose}
\label{sec:docker-compose}

\subsection{Composing the API and Its Supporting Services}

\subsection{Volumes and Networks}

\subsection{Environment Configuration and Secrets}

%==============================================================================
\section{Model Serving at Scale}
\label{sec:serving}

\subsection{Inference Engines: \texttt{vLLM} and \texttt{llama.cpp}}

\subsection{Batching}

\subsection{The Latency--Throughput Tradeoff}

%==============================================================================
\section{Cloud Deployment}
\label{sec:cloud}

\subsection{Compute and Storage: EC2 and S3}

\subsection{Managed Training and Endpoints: SageMaker}

%==============================================================================
\section{Monitoring}
\label{sec:monitoring}

\subsection{Structured Logging}

\subsection{Drift Detection}

\subsection{The Retraining Trigger}
